\documentclass[
11pt,notheorems,hyperref={pdfauthor=whatever}
]{beamer}

\input{loadslides.tex} % Loads packages and some defined commands

\title[
% Text entered here will appear in the bottom middle
]{Réunion du 12/05/2026}

\subtitle{IFPEN/INRIA}

\author[
% Text entered here will appear in the bottom left corner
]{
    Arthur Timbeau, 
    Adam Rozentalis,
    Julien Salomon
}


\AtBeginSection[]{
  \begin{frame}[plain]
    \vfill
    \centering
    \usebeamerfont{title}\insertsection
    \vfill
  \end{frame}
}

\begin{document}

% Generate title page

\setbeamertemplate{footline}{}
\begin{frame}
  \titlepage
\end{frame}

\setbeamertemplate{footline}{
  \hfill\insertframenumber/\inserttotalframenumber\hspace{1em}\vspace{0.5em}
}

\begin{frame}
    \tableofcontents
\end{frame}

\section{Introduction}

\subsection{Brève introduction de la présentation}
\begin{frame}{Brève introduction de la présentation}
    Les modèles dit de la \textcolor{blue}{\textit{branche 1}} sont les modèles qui ont 
    été développé en premier et présentéé lors de la réunion du 22/04/26.
    
    \vspace{0.3cm}

    Ceux de la \textcolor{blue}{\textit{branche 2}} sont ceux qui ont été écrit tout récemment. 
\end{frame}

\subsection{Contenu du document}
\begin{frame}{Contenu du document}
    \begin{itemize}
        \item Rappels sur les travaux précédemment réalisé et présentation des stratégies 
        \item Présentation des corrections sur les tests de la branche 1 et résumé des performances
        \item Introduction de la branche 2 et comparaison avec la branche 1
        \item Présentation des résultats numériques
    \end{itemize}    
\end{frame}

\section{Rappels sur les stratégies employées}

\subsection{Corrections locales/Globales}

\begin{frame}{Deux familles de correction}
    Quatres familles de modèles sont envisagées :
    \begin{center}
    \begin{itemize}
    \item Les corrections \textbf{locales} (L) : les \textcolor{red}{rayons} et les \textcolor{blue}{azimuts} sont des paramètres d’entrés
    \item Les corrections à rayon fixé ($G_r$ ) : les \textcolor{blue}{azimuts}  sont des paramètres d’entrés.
    \item Les corrections à azimut fixé ($G_a$ ) : les \textcolor{red}{rayons} sont des paramètres d’entrés
    \item Les corrections \textbf{globales} (G) : les \textcolor{red}{rayons} et les \textcolor{blue}{azimuts}  sont fixés
    \end{itemize}
\end{center}
\end{frame}


\begin{frame}{Approche L}
    Approche locale : inférences sur les efforts par sections radiale. 

    \vspace{0.3cm}

    \begin{columns}[c]
        \column{0.44\textwidth}
        \textbf{Inputs :}
        \begin{itemize}\small
            \item Position radiale $r \in \left[R_{hub}, R\right]$
            \item Position azimutale $\theta \in \left[0,2\pi\right]$
            \item Vitesse de rotation $\Omega \in \mathbb{R}$
            \item Vitesse freestream $U_{\infty} \in \mathbb{R}$
            \item Yaw $\gamma \in \left[-\frac{\pi}{4}, \frac{\pi}{4}\right]$
        \end{itemize}

        \column{0.10\textwidth}
        \centering
        
\begin{tikzpicture}
            \draw[-stealth, line width=4pt, color=blue!70]
                (0, 0) -- (0.9, 0);
        \end{tikzpicture}

        \column{0.44\textwidth}
        \textbf{Outputs :}
        \begin{itemize}\small
            \item Force normale $F_N \in \mathbb{R}$
            \item Force tangentielle $F_T \in \mathbb{R}$
        \end{itemize}
    \end{columns}
\end{frame}

\begin{frame}{Approche $G_r$}
    Approche à azimut fixé : inférences sur les efforts en fonction d’un angle azimutal sur une grille de p
    éléments.

    \vspace{0.3cm}

    \begin{columns}[c]
        \column{0.44\textwidth}
        \textbf{Inputs :}
        \begin{itemize}\small
            \item Position radiale $r \in \left[R_{hub}, R\right]$
            \item Vitesse de rotation $\Omega \in \mathbb{R}$
            \item Vitesse freestream $U_{\infty} \in \mathbb{R}$
            \item Yaw $\gamma \in \left[-\frac{\pi}{4}, \frac{\pi}{4}\right]$
        \end{itemize}

        \column{0.10\textwidth}
        \centering
        
\begin{tikzpicture}
            \draw[-stealth, line width=4pt, color=blue!70]
                (0, 0) -- (0.9, 0);
        \end{tikzpicture}

        \column{0.44\textwidth}
        \textbf{Outputs :}
        \begin{itemize}\small
            \item Force normale $F_N \in \mathbb{R}^p$
            \item Force tangentielle $F_T \in \mathbb{R}^p$
        \end{itemize}
    \end{columns}
\end{frame}

\begin{frame}{Approche $G_a$}
    Approche à azimut fixé : inférences sur les efforts en fonction d’un rayon sur une grille de n éléments.

    \vspace{0.3cm}

    \begin{columns}[c]
        \column{0.44\textwidth}
        \textbf{Inputs :}
        \begin{itemize}\small
            \item Position azimutale $\theta \in \left[0,2\pi\right]$
            \item Vitesse de rotation $\Omega \in \mathbb{R}$
            \item Vitesse freestream $U_{\infty} \in \mathbb{R}$
            \item Yaw $\gamma \in \left[-\frac{\pi}{4}, \frac{\pi}{4}\right]$
        \end{itemize}

        \column{0.10\textwidth}
        \centering
        
\begin{tikzpicture}
            \draw[-stealth, line width=4pt, color=blue!70]
                (0, 0) -- (0.9, 0);
        \end{tikzpicture}

        \column{0.44\textwidth}
        \textbf{Outputs :}
        \begin{itemize}\small
            \item Force normale $F_N \in \mathbb{R}^n$
            \item Force tangentielle $F_T \in \mathbb{R}^n$
        \end{itemize}
    \end{columns}
\end{frame}
\begin{frame}{Approche G}
Approche globale : inférences sur les efforts sur toute les sections de la pale et tous les azimuts.

    \vspace{0.3cm}

    \begin{columns}[c]
        \column{0.44\textwidth}
        \textbf{Inputs :}
        \begin{itemize}\small
            \item Vitesse de rotation $\Omega \in \mathbb{R}$
            \item Vitesse freestream $U_{\infty} \in \mathbb{R}$
            \item Yaw $\gamma \in \left[-\frac{\pi}{4}, \frac{\pi}{4}\right]$
        \end{itemize}

        \column{0.10\textwidth}
        \centering
        
\begin{tikzpicture}
            \draw[-stealth, line width=4pt, color=blue!70]
                (0, 0) -- (0.9, 0);
        \end{tikzpicture}

        \column{0.44\textwidth}
        \textbf{Outputs :}
        \begin{itemize}\small
            \item Force normale $F_N \in \mathbb{R}^{n \times p}$
            \item Force tangentielle $F_T \in \mathbb{R}^{n \times p}$
        \end{itemize}
    \end{columns}    
\end{frame}


\begin{frame}{Résumé}
    Les grandes familles de correction 
    \begin{itemize}
        \item \centering  $L$ :     inputs = $(r, \theta, \Lambda)$     || outputs = $(F_n, F_t) \in \mathbb{R}^2$
        \item \centering  $G_r$ :   inputs = $(r, \Lambda)$             || outputs = $(\mathbf{F_n}, \mathbf{F_t}) \in \mathbb{R}^{2n}$
        \item \centering  $G_a$ :   inputs = $(\theta, \Lambda)$        || outputs = $(\mathbf{F_n}, \mathbf{F_t}) \in \mathbb{R}^{2p}$
        \item \centering  $G$ :     inputs = $(\Lambda)$                || outputs = $(\mathbf{F_n}, \mathbf{F_t}) \in \mathbb{R}^{2np}$
    \end{itemize}

     \vspace{0.3cm}

     Où $\Lambda$ est un vecteur contenant d'autres paramètres indépendant de la géométrie et de la discrétisation.
\end{frame}

\subsection{Approche directe/intermédiaire}

\begin{frame}{Approche directe}
    \textbf{Approche directe} : corriger directement les efforts retournés par la BEM : 
    
    \vspace{0.4cm}
    \begin{center}
    \begin{tikzpicture}[node distance=1.5cm]
        \node[draw, rounded corners, fill=blue!10, minimum width=3cm, minimum height=1.1cm, align=center]
            (bem) {Efforts BEM : \\ $\bigl(F_n^{BEM},\, F_t^{BEM}\bigr)$};
        \node[draw, fill=orange!20, minimum width=2cm, minimum height=1.1cm, align=center, right=1.5cm of bem]
            (nn)  {NN\\SkyNet};
        \node[draw, rounded corners, fill=green!15, minimum width=3cm, minimum height=1.1cm, align=center, right=1.5cm of nn]
            (skn) {Sorties inférées : \\ $\bigl(F_n^{SKN},\, F_t^{SKN}\bigr)$};
        \draw[-stealth, line width=2pt] (bem) -- (nn);
        \draw[-stealth, line width=2pt] (nn)  -- (skn);
    \end{tikzpicture}
    \end{center}

    Dans ce cas, le vecteur $\Lambda$ contient les efforts normaux calculés par la BEM (soit une valeur, soit un distribution radiale ou azimutale ou soit 
    une matrice de $n\times p$ coefficients.)
\end{frame}

\begin{frame}{Approche intermédiaire}
    Approche intermédiaire : considérer une variable en amont des efforts (présente dans BEM et FWV) et inférer dessus. 

        \begin{center}
    \begin{tikzpicture}[node distance=1.5cm]
        \node[draw, rounded corners, fill=blue!10, minimum width=3cm, minimum height=1.1cm, align=center]
            (bem) {Sorties BEM : \\ $X_{BEM}$};
        \node[draw, rounded corners, fill=orange!15, minimum width=3cm, minimum height=1.1cm, align=center, right=1.5cm of bem]
            (var_inter) {Sorties inférées : \\ $\text{SkyNet}(X_{BEM}) = Y_{SKN}$};

        \draw[-stealth, line width=2pt] (bem) -- (var_inter);
    \end{tikzpicture}

    \vspace{0.3cm}
    Où $G$ représente les différentes formules permettant de calculer les efforts à partir de la variable inférée.

    \vspace{0.2cm}

    Par exemple, les deux modèles utilisent dans le calcul des efforts la \textbf{vitesse effective} couplé à \textbf{l'angle d'attaque}.
    \begin{columns}[c]
        \column{0.44\textwidth}
        \textbf{BEM :}
        Au niveau de la section de rayon $r$
        \begin{itemize}\small
            \item vitesse effective : $V_{eff} = \sqrt{\left[ U_{\infty}(1-a)^2 + ^\Omega r(1+a')^2\right]}$
            \item angle d'attaque : $\alpha = \varphi -\left(\beta + \theta_{pitch}\right)$
        \end{itemize}

        \column{0.10\textwidth}
        \centering
        
\begin{tikzpicture}
            \draw[-stealth, line width=4pt, color=blue!70]
                (0, 0) -- (0.9, 0);
        \end{tikzpicture}

        \column{0.44\textwidth}
        \textbf{FWV :}
        Au niveau de la section de rayon $r$ 
        \begin{itemize}\small
            \item Biot-Savart :  $V_{eff} = \frac{\Gamma(||r_1|| + ||r_2||)(r_1 \times r_2)}{4\pi ||r_1|| \, ||r_2||(||r_1||\,||r_2|| + (r_1 | r_2))}$
            \item angle d'attaque :  $\alpha = \text{arctan}\bigl(\frac{(u|N)}{(u|T)}\bigr)$
        \end{itemize}
    \end{columns}    
    \end{center}
\end{frame}

\begin{frame}{Variables d'entrée}
    Que ce soit pour la variante intermédiaire ou directe, les variables d'entrées de ces modèles ont des entrées BEM 
    \begin{itemize}
        \item Sous la forme de scalaires dans l'approche $L$ (soit $ \Lambda \in \mathbb{R}^2$)
        \item Sous la forme de vecteurs de $n$ coefficients dans l'approche $G_r$ (soit $ \Lambda \in \mathbb{R}^{2n}$)
        \item Sous la forme de vecteurs de $p$ coefficients dans l'approche $G_a$ (soit $ \Lambda \in \mathbb{R}^{2p}$)
        \item Sous la forme d'une matrice de $n \times p$ coefficients dans l'approche $G$ (soit $ \Lambda \in \mathbb{R}^{2np}$)
    \end{itemize}
\end{frame}

\begin{frame}{Stratégies sans BEM : créer un challenger}
    \begin{center}
    En parallèle, des modèles émulant directement FWV (sans utiliser d'autre modèle) sont écrits dans le but de challenger les modèles corrigeant la BEM.\\ 
    Objectif : \textbf{Montrer qu'un réseau de neurones informés par un modèle physique performe mieux qu'un modèle non informé}.
    \end{center}
\end{frame}

\begin{frame}{Ensemble des stratégies possible}
    \begin{center}
        L'ensemble des stratégies possibles se résume par la formule suivante : 
        \begin{align*}
            \bigl(\text{local}, \text{global}\bigr) \times \bigl(\text{Avec BEM}, \text{Sans BEM}\bigr) \times \bigl(\text{Variable intermédiaire}, \text{Directe}\bigr)
        \end{align*}
    \end{center}
\end{frame}


\subsection{Architure du réseau de neurones}

\begin{frame}{Réseau de neurones}
    Le réseau de neurones est écrit avec la librairie PyTorch. Les modèles sont des \textit{Multi Layer Perceptron} (MLP). Les couches cachées sont de même taille. 
    
    \vspace{0.3cm}
    
    La couche de d'entrée est de dimension (\textbf{variable d'entrée }$\times$ \textbf{nombre de neurones}) et 

    la couche de sortie est de dimension (\textbf{nombre de neurones} $\times$  \textbf{variable de sortie}).

    Les fonctions d'activations sont les mêmes pour chaque couche : il s'agit de la \textit{REctified Linear Unit} (RELU) : 
    \begin{align*}
        \text{RELU}(x) = \text{max}(0,x)        
    \end{align*}
\end{frame}

\begin{frame}{Hyperparamètres}
    En plus des hyperparamètres classiques, un \textit{taux de dropout} est utilisé afin de s'assurer que chaque neurone est mobilisé lors de l'apprentissage.
    
    \vspace{0.2cm}

    Les hyperparamètres sont choisit avec la librairie \textit{optuna}. Plusieurs entrainements sont lancés avec différents hyperparamètres choisit dans la plage
    suivante : 
    
    \begin{itemize}
        \item nombre de couches cachées : entre $2 \text{ et } 8$ (chosit entre 2 et 6)
        \item nombre de neurones : entre $64$ et $512$ (chosit entre 277 et 493 )
        \item taux de dropout : entre $0$ et $0.5$ (chosit entre $4.952e^{-5}$ et $0.0313$)
        \item taux d'apprentissage : entre $1e^{-4} \text{ et } 1e^{-2}$ (choisit entre $0.0001$ et $0.0061$)
    \end{itemize}
    
    Ils sont choisit comme étant ceux minimisant une erreur de validation croisée à 3 folds.
\end{frame}

\begin{frame}{Entrainement}
    Pour accélerer l'entrainement, une méthode de \textit{Mini-Batch} à été mise en place. L'entrainement des modèles de la famille 
    $G$ n'a pas nécéssité de batch. Pour $L$, les batchs contiennent $1024$ éléments (taille classiques), tandis que pour les familles $G_A \text{ et } G_R$
    on nécéssité un réajustement à 52 éléments par batchs.

    \vspace{0.3cm}

    On notera que pour un même nombre d'epoch (1000), le temps d'entraînement par epoch croît : on compte 2 secondes pour les modèle $G$, 20 secondes pour les modèles $G_R$
    et 2 minutes pour les modèles $L$.
    
\end{frame}

%%%%%%%%%%%%%%%%%%%%%%%%%%%%%%%%%%%%%%%%%%%%%%%%%%%%%%%%%%%%%%%%%%%%%%%%%%%%%%%%%%%%%%%%%%%%%%%%%%%%%%%%%%%%%%%%%%%%%%%%%%%%%
%%%%%                                   BRANCHE 1                                                                       %%%%% 
%%%%%%%%%%%%%%%%%%%%%%%%%%%%%%%%%%%%%%%%%%%%%%%%%%%%%%%%%%%%%%%%%%%%%%%%%%%%%%%%%%%%%%%%%%%%%%%%%%%%%%%%%%%%%%%%%%%%%%%%%%%%%

\section{Branche 1}

\subsection{Retour sur les benchmarks des solveurs}
\begin{frame}{Validation erronée}
    Les tests montré lors de la dernière réunion présentaient plusieurs problèmes. 
    \begin{itemize}
        \item (I) Les géométries utilisées pour faire tourner les solveurs n'étaient pas les mêmes         
        \item (II) Certains calculs étaient faux.
    \end{itemize}
    \vspace{0.2cm}
    Toutes ces incohérences ont été corrigées. La conclusion est celle attendue : 
    
    \vspace{0.3cm}

    \begin{align}
            \textbf{SVEN} \geq \textbf{BEM avec yaw} \text{ (en terme de fidéflité à CASTOR)}
    \end{align}
\end{frame}

\begin{frame}{Validation corrigée}
    \begin{figure}[h]
    \centering 
        \includegraphics[width=0.497\textwidth]{sanitycheck_Yaw15°_forces_rayons.png}
        \includegraphics[width=0.497\textwidth]{wrong_val.png}
        \caption{Validation corrigée à gauche et erronée à droite} 
        \label{fig:val} 
    \end{figure}
\end{frame}

\begin{frame}{Mesure de fidélité}
    \begin{columns}[c]
        \column{0.44\textwidth}
            \captionof{table}{Mesures de fidélité à CASTOR pour $F_n$}
            \begin{tabular}{c c c}
                \hline
                Modèle               & BEM   & SVEN  \\
                \hline
                RMSE (stricte) (N/m) & 26.82 & 13.20 \\
                RMSE (rel.) (\%)     & 12.9\%& 6.4\% \\
                ERR Résultante (\%)  & 9.1\% & 0.1\% \\
                Wassertstein         & 18.806& 3.391 \\
                Corrélation (\%)     & 99.4\% & 99.6\% \\
                \hline
            \end{tabular}
        
        \column{0.44\textwidth}
            \captionof{table}{Mesures de fidélité à CASTOR pour $F_t$}
            \begin{tabular}{c c c}
                \hline
                Modèle               & BEM   & SVEN     \\
                \hline
                RMSE (stricte) (N/m) & 4.82   & 2.98    \\
                RMSE (rel.) (\%)     & 31.2\% & 19.3\%  \\
                ERR Résultante (\%)  & 22.3\% & 1.4\%   \\
                Wassertstein         & 3.482  & 0.482   \\
                Corrélation (\%)     & 96.2\% & 97.2\%  \\
                \hline
            \end{tabular}
    \end{columns}    
\end{frame}

\begin{frame}{Nouvelle mesure de validation}
    Une nouvelle méthode de mesure de qualité d'une prédiction mise en place consiste à regarder les rendements calculés à partir des informations inférées. 
    Pour valider cette mesure, elle a été utilisé pour tester la fidélité de BEM et SVEN à CASTOR.

    \begin{figure}[h]
        \centering 
        \includegraphics[width=0.7\textwidth]{sanitycheck_rendement.png}
        \caption{Rendements} 
        \label{fig:rend_fid} 
    \end{figure}
    
\end{frame}

\begin{frame}{Élimination de la reconstruction BEM à partir de FVW}
    L'une des stratégies consistait à corriger les sorties $\big(a,a',\varphi\big)$ de la BEM par rapport à leurs analogues reconstruit à partir de 
    FVW. 

    \vspace{0.5cm}
    
    \begin{columns}[c]
        \column{0.44\textwidth}
        \textbf{BEM} (simplifiée) :
        Au niveau de la section de rayon $r$
        \begin{itemize}
            \item $\tan(\varphi) = \frac{1-a}{\lambda_r(1+a')}$
            \item $\frac{a}{1-a} = \frac{\sigma_{\lambda_r}}{4F(\varphi)\sin(\varphi)^2}C_L(\varphi - \beta_{r})cos(\varphi)$ 
        \end{itemize}
            
        \column{0.10\textwidth}
        \centering
        
\begin{tikzpicture}
            \draw[-stealth, line width=4pt, color=blue!70]
                (0, 0) -- (0.9, 0);
        \end{tikzpicture}
        
\begin{tikzpicture}
            \draw[-stealth, line width=4pt, color=blue!70]
                (0.9, 0) -- (0, 0);
        \end{tikzpicture}

        \column{0.44\textwidth}
        \textbf{FWV :}
        Au niveau de la section de rayon $r$ 
        \begin{itemize}\small
            \item $ \varphi = \alpha + \beta_r + \theta_{pitch}$
            \item $a = 1 - \sin(\varphi)\frac{V_{eff}}{U_{\infty}}$
        \end{itemize}
    \end{columns}    
    
    \vspace{0.5cm}
    
    Cette approche à été mise de côté pour plusieurs raisons : 
    \begin{itemize}
        \item Le calcul BEM de $a$ est plus sophistiqué que celui de FVW
        \item Adaptation du calcul complexe : problème d'unicité dans les valeurs retournées
    \end{itemize}

    \vspace{0.3cm}
    
    Elle sera peut-être réexplorée après une étude plus approfondie.
\end{frame}

\subsection{Conclusion sur la branche 1}

\begin{frame}{Stratégies testés}
    Stratégies testées sur la branche 1 : 
    \[
        \bigl(\text{local}, \text{GA}, \text{GR}, \text{global}\bigr) \times \bigl(\text{Avec BEM}, \text{Sans BEM}\bigr) \times \bigl((V_{eff},\alpha) , \text{Efforts normaux et tangentiels}\bigr)
    \]
    
\end{frame}

\begin{frame}{Résultats bruts}
    \captionof{table}{\textbf{Résultats bruts pour la branche 1}}
    \begin{tabular}{c c c c c c c}
        \hline
        Modèle & Epoch & Score Global & RMSE $F_n$ & RMSE $F_t$ & Wass $F_n$ & Wass $F_t$ \\
        \hline
        GR\_0\_f & 1000 & 0.30\% & 0.07\% & 0.23\% & 0.085 & 0.023 \\
        GR\_0\_v & 1000 & 0.32\% & 0.08\% & 0.24\% & 0.110 & 0.022 \\
        GR\_1\_f & 1000 & 0.70\% & 0.07\% & 0.63\% & 0.100 & 0.036 \\
        GR\_1\_v & 1000 & 0.65\% & 0.24\% & 0.41\% & 0.370 & 0.050 \\
        GR\_2\_f & 1000 & 0.42\% & 0.10\% & 0.32\% & 0.129 & 0.031 \\
        GR\_2\_v & 1000 & 0.39\% & 0.10\% & 0.29\% & 0.137 & 0.029 \\
        \hline
        L\_0\_f  & 1000 & 23.02\% & 8.62\% & 14.40\% & 4.218 & 0.567 \\
        L\_0\_v  & 1000 & 9.96\%  & 2.29\% & 7.67\%  & 2.167 & 0.458 \\
        L\_1\_f  & 1000 & 1.81\%  & 0.27\% & 1.54\%  & 0.387 & 0.115 \\
        L\_1\_v  & 1000 & 3.34\%  & 0.88\% & 2.46\%  & 1.109 & 0.241 \\
        L\_2\_f  & 1000 & 0.94\%  & 0.19\% & 0.76\%  & 0.238 & 0.055 \\
        L\_2\_v  & 1000 & 3.39\%  & 0.87\% & 2.52\%  & 0.921 & 0.165 \\
        \hline
        GA\_0\_f & 1000 & 40.61\% & 11.75\% & 28.86\% & 12.112 & 1.712 \\
        GA\_0\_v & 1000 & 18.25\% & 3.86\%  & 14.39\% & 6.525  & 0.951 \\
        GA\_1\_f & 1000 & 11.25\% & 2.96\%  & 8.29\%  & 4.379  & 0.832 \\
        GA\_1\_v & 1000 & 34.00\% & 8.30\%  & 25.71\% & 10.598 & 1.719 \\
        GA\_2\_f & 1000 & 16.23\% & 4.20\%  & 12.04\% & 5.256  & 0.968 \\
        GA\_2\_v & 1000 & 31.41\% & 7.86\%  & 23.55\% & 9.997  & 1.860 \\
        \hline
    \end{tabular}
\end{frame}

\begin{frame}{Résumé des résultats bruts}
    L'utilisation de la BEM pour les modèles donne de meilleurs résultats dans l'approche locale.
    
    \vspace{0.6cm}
    
    Pour les stratégies semi-locale, l'information de la BEM les rends moins efficaces. On note aussi que l'approche à rayon fixé (GA) est très peu efficace. 

    \vspace{0.6cm}

    Ces résultats sont plutôt contre-intuitif : on s'attendrait à ce qu'un modèle informé ,au travers de plus de variable, par une physique plus simple
    devrait permettre de meilleures prédiction sur la physique plus sophistiquée. 
\end{frame}


\subsection{Graphiques}

\begin{frame}{Graphiques : $L$}
    \begin{figure}[h]
    \centering 
        \includegraphics[width=0.4\textwidth]{branche1_summary_Meilleur_L.png}
        \includegraphics[width=0.4\textwidth]{branche1_summary_Pire_L.png}
        \caption{Meilleures performances à gauche et pire à droite} 
        \label{fig:L_perf_1} 
    \end{figure}
\end{frame}

\begin{frame}{Graphique : $G_R$}
        \begin{figure}[h]
        \centering 
        \includegraphics[width=0.4\textwidth]{branche1_summary_Meilleur_GR.png}
        \includegraphics[width=0.4\textwidth]{branche1_summary_Pire_GR.png}
        \caption{Meilleures performances à gauche et pire à droite} 
        \label{fig:GR_perf_1} 
    \end{figure}
\end{frame}

\begin{frame}{Graphique : $G_A$}
        \begin{figure}[h]
        \centering 
        \includegraphics[width=0.4\textwidth]{branche1_summary_Meilleur_GA.png}
        \includegraphics[width=0.4\textwidth]{branche1_summary_Pire_GA.png}
        \caption{Meilleures performances à gauche et pire à droite} 
        \label{fig:GA_perf_1} 
    \end{figure}
\end{frame}

%%%%%%%%%%%%%%%%%%%%%%%%%%%%%%%%%%%%%%%%%%%%%%%%%%%%%%%%%%%%%%%%%%%%%%%%%%%%%%%%%%%%%%%%%%%%%%%%%%%%%%%%%%%%%%%%%%%%%%%%%%%%%
%%%%%                                   BRANCHE 2                                                                       %%%%% 
%%%%%%%%%%%%%%%%%%%%%%%%%%%%%%%%%%%%%%%%%%%%%%%%%%%%%%%%%%%%%%%%%%%%%%%%%%%%%%%%%%%%%%%%%%%%%%%%%%%%%%%%%%%%%%%%%%%%%%%%%%%%%


\section{Branche 2}

\subsection{Présentation}

\begin{frame}{Présentation}
    Les modèles de la branche 2 ont pour but de permettre l'implémentation de la la stratégie globale (abrégée G). 
    
    \vspace{0.3cm}

    L'idée consiste à travailler avec une grille d'azimutale de $p \in \mathbb{N}$ éléments et une grille radiale de $n \in \mathbb{N}$ éléments. 
    Dans cette stratégie, les variables sur lesquels le modèle peut s'entraîner sont le yaw et le TSR. Les simulations effectuées avaient $p = 36$ et $n = 36$. 

    \vspace{0.3cm}

    L'objectif à terme est de constituer un dataset contenant les diverses variables d'intérêts pour plusieurs combinaison de yaw et de TSR. Par soucis
    de simplicité, le TSR à été fixé pour pouvoir commencer les 1er tests.

    \vspace{0.3cm}

    Les codes ont été écrit de sorte que les autres stratégies ait pu être essayées.
\end{frame}

\subsection{Comparaison conceptuelle avec la branche 1}
\begin{frame}{Comparaison conceptuelle avec la branche 1}
La différence majeure entre les modèles de la branche 1 et la branche 2 consiste en les \textbf{paramètres d'entrées}. 

\vspace{1cm}

\begin{columns}[c]
        \column{0.44\textwidth}
        Variables d'entrées \textbf{Branche 1} :  
        \begin{itemize}\small
            \item Informations BEM (forces ou variable intermédiaire)
            \item $r$ ou $\theta$ selon le choix de la famille de modèle (L,G,GA,GR)
        \end{itemize}

        \column{0.10\textwidth}
        \centering

        \column{0.44\textwidth}
        Variables d'entrées \textbf{Branche 2} :
        \begin{itemize}\small
            \item Informations BEM (forces ou variable intermédiaire)
            \item $r$ ou $\theta$ selon le choix de la famille de modèle (L,G,GA,GR)
            \item Yaw $\gamma \in \left[-\frac{\pi}{4}, \frac{\pi}{4}\right]$
            \item TSR $\in \left[4, 12\right] $
        \end{itemize}
    \end{columns}

    \vspace{1cm}

    Pour deux modèles d'une même famille, celui de la branche 1 est a priori moins "généralisable" que celui de la branche 2. Ce dernier est capable d'inférer sur des configurations métérologiques 
    différentes.  
\end{frame}

\subsection{Performances de la branche 2}
\begin{frame}{Résultats brut}

    \captionof{table}{\textbf{Résumé brut pour la branche 2}}
    \begin{tabular}{c c c c c c c}
        \hline
        Modèle & Epoch & Score Global & RMSE $F_n$ & RMSE $F_t$ & Wass $F_n$ & Wass $F_t$ \\
        \hline
        GR\_0\_f & 1000 & 0.88\% & 0.15\% & 0.72\% & 0.137 & 0.037 \\
        GR\_1\_f & 1000 & 1.33\% & 0.31\% & 1.02\% & 0.202 & 0.041 \\
        \hline
        GR\_0\_v & 1000 & 4.08\% & 1.09\% & 2.99\% & 0.783 & 0.167 \\
        GR\_1\_v & 1000 & 1.04\% & 0.12\% & 0.92\% & 0.103 & 0.016 \\
        \hline
        L\_0\_f  & 1000 & 14.10\% & 5.55\% & 8.55\% & 4.242 & 0.315 \\
        L\_1\_f  & 1000 & 1.61\% & 0.33\% & 1.28\% & 0.265 & 0.065 \\
        \hline
        L\_0\_v  & 1000 & 11.52\% & 2.58\% & 8.95\% & 1.844 & 0.498 \\
        L\_1\_v  & 1000 & 4.22\% & 1.20\% & 3.02\% & 0.904 & 0.128 \\
        \hline
        G\_0\_f  & 1000 & 3.23\% & 0.69\% & 2.54\% & 0.658 & 0.160 \\
        G\_1\_f  & 1000 & 3.95\% & 0.81\% & 3.14\% & 0.744 & 0.118 \\
        \hline
        G\_0\_v  & 1000 & 3.81\% & 0.97\% & 2.84\% & 0.853 & 0.183 \\
        G\_1\_v  & 1000 & 5.96\% & 1.95\% & 4.02\% & 0.962 & 0.180 \\
        \hline
    \end{tabular}
\end{frame}

\begin{frame}{Résumé des résutats bruts}
    Les mesures sont, dans ces cas, toutes équivalentes les unes par rapport aux autres.

    \vspace{0.3cm}
    
    Les modèles de la famille $L$ semblent beaucoup plus performant lorsqu'ils ont l'information de la BEM.

    \vspace{0.3cm}

    Dans les autres cas, l'approche $(V_{eff}, \alpha)$ semble meilleur avec que sans la BEM. L'approche directe dans ces cas est meilleure sans BEM.
\end{frame}

\begin{frame}{Graphiques : $L$}
        \begin{figure}[h]
        \centering 
        \includegraphics[width=1\textwidth]{branche2_synthese_famille_L.png}
        \caption{Plot des rendements} 
        \label{fig:rend_L_2} 
    \end{figure}
\end{frame}

\begin{frame}{Graphiques : $G_R$}
    \begin{figure}[h]
        \centering 
        \includegraphics[width=1\textwidth]{branche_synthese_famille_GR.png}
        \caption{Plot des rendements} 
        \label{fig:rend_GR_2} 
    \end{figure}
\end{frame}

\begin{frame}{Graphiques : $G$}
        \begin{figure}[h]
        \centering 
        \includegraphics[width=1\textwidth]{branche2_synthese_famille_G.png}
        \caption{Plot des rendements} 
        \label{fig:rend_G_2} 
    \end{figure}
\end{frame}


\subsection{Comparaison des performances entre la branche 1 et 2 }

\begin{frame}{Comparaison des résultats}
L'ajout du yaw comme paramètre d'entrée des modèles fait baisser leur précision dans l'ensemble au sens où la plus faible erreur de précision est plus grande lorsque
le yaw est un paramètre. En revanche, elle est "beaucoup plus stable", c'est-à-dire que la plus grande erreur de précision est plus faible lorsque le yaw est un paramètre pour la mesure de Wassertstein.

\vspace{0.3cm}

L'ajout du yaw donne plus de poid aux modèles informés par la BEM : un modèle informé par une physique "grossière" permet de mieux prédire les sorties d'une physique plus "fine". 


\end{frame} 
\end{document}